% Library unit: the long side (rook endings). Printed once, before the first position that uses it; see lib/README.md.
\paragraph{Long side}\label{long-side}\index{Rook!Long side}

Every pawn except a rook's pawn divides the board into two parts of different width. The wider
part is the
\term{long-side}.
In \xref{diagram}{side-checks} the e-pawn has four files on its left (a, b, c, d): that is
the long side.

Because \term{side-checks} need room, the defending rook belongs on the long
side. Put Black's rook on h2 instead of a2 in \xref{diagram}{side-checks}: it is too close to
the pawn to give side checks, and White wins.
